%% 自動生成（phys/gen_phys.py）
\phtocgrid{\phgi{保存則と相似の地図}{4}\phgi{この本の使い方}{6}}
{\color{bkLine}\hrule height 0.35pt}\vspace{2.4mm}
\phtocchap{原-0}{スマホと原子炉、どちらも原子の物理で動く？}{原子1個は、どれくらい小さい？}{}{3}{3}{8}{J}
\phtocpart{第1部}{飛び出す：光と電子の二重性}{12}{1}
\phtocchap{原-1}{原子の中をのぞく：電子と原子核の発見}{原子の中は、ほとんどからっぽだって本当？}{2}{4}{4}{14}{C}
\phtocchap{原-2}{光はつぶでもある：光電効果と光子}{熱した鉄は、なぜ温度で光の色が変わる？}{5,7}{4}{4}{30}{C}
\phtocchap{原-3}{電子は波でもある：ド・ブロイ波と不確定性}{電子顕微鏡は、なぜふつうの顕微鏡より小さなものが見える？}{4,8,9}{4}{4}{46}{C}
\phtocchap{原-4}{原子が出す光：エネルギー準位とスペクトル}{花火の色は、なぜ材料によって変わる？}{4}{4}{4}{62}{C}
\phtocpart{第2部}{すりぬける：量子力学の入口}{78}{2}
\phtocchap{原-5}{波で書く電子：シュレディンガー方程式}{電子は、原子の中のどこにいる？}{6,7,9}{5}{5}{80}{C}
\phtocchap{原-6}{壁をすりぬける：トンネル効果と調和振動子}{USBメモリは、どうやって電子を閉じこめて記録する？}{5,6}{5}{5}{98}{C}
\phtocchap{原-7}{電子の席の決まり：原子と周期表}{ネオンは、なぜほかの物質とほとんど結びつかない？}{4,8}{5}{5}{116}{C}
\phtocbreak
\phtocpart{第3部}{切りかえる：固体と半導体}{134}{3}
\phtocchap{原-8}{固体の中の電子：エネルギーバンド}{金属は電気を通すのに、ガラスはなぜ通さない？}{3,4,5}{5}{5}{136}{C}
\phtocchap{原-9}{一方通行の電気：pn接合とLED・太陽電池}{LEDは、どうやって電気を直接、光に変える？}{3,5}{5}{5}{154}{C}
\phtocchap{原-10}{スイッチと増幅：トランジスタと電子回路}{スマホのチップは、どうやって計算している？}{2,7}{4}{5}{172}{C}
\phtocpart{第4部}{分かれる：原子核と放射線}{190}{4}
\phtocchap{原-11}{質量はエネルギー：相対論と原子核}{太陽は、何をエネルギーのもとにして光っている？}{2,4,5}{4}{5}{192}{C}
\phtocchap{原-12}{こわれる原子核：放射性崩壊と半減期}{大昔の木や骨の年代は、どうやってわかる？}{3,6}{4}{5}{210}{C}
\phtocchap{原-13}{放射線が物にあたる：放射線と物質}{X線写真で、なぜ骨が白くうつる？}{3,5}{5}{5}{228}{C}
\phtocchap{原-14}{放射線をはかる・身を守る：計測と防護}{目に見えない放射線を、どうやってはかる？}{3,6,8}{4}{4}{246}{C}
\phtocchap{原-15}{原子炉の中の中性子：核分裂と原子炉}{原子炉は、どうやって反応を一定に保つ？}{6,8,9}{5}{5}{262}{C}
\phtocchap{原-16}{原子の力とつきあう}{使い終わった核燃料は、どうする？}{}{4}{4}{280}{R}
\par\vspace{1.5mm}
\phtocgrid{\phgi{こたえ}{288}\phgi{記号カード}{296}\phgi{核種のデータ・物性値・定数}{299}\phgi{\mbox{なぜ？}の道すじ}{304}\phgi{できるようになったことチェック表}{306}\phgi{つぎの巻へ}{308}\phgi{さくいん}{310}}
